\subsection{Why the Generated State Represents the Contract}
\label{sec:contract-adequacy}
A \emph{local update history} $\ell=z(a_1,o_1)\cdots(a_n,o_n)$ begins at an old closed-loop tuple and records each event and its realized local successors. Components follow their active-version transitions and local transfers; boundaries preserve physical states. Monitor memories fold their active intervals: stop excludes its own boundary; start installs its initializer after the boundary. Initially inherit old memories, initialize interval monitors and release old control, leaving new monitors inactive.

Ordinary events retain all combinations of participating local outcomes. Commands require UC quiescence, unexecuted status, completed predecessors and a transfer/initializer value where applicable. Quiescence includes unsafe UC outcomes. A source policy selects controllable labels and must retain every UC and every permitted outcome. Every maximal history must safely reach a quiescent matching target after all commands. Technical Appendix~\ref{ta:source_synthesis} gives the full recurrences.

\begin{proposition}[Adequacy of the generated update game]
\label{prop:contract-adequacy}
For an eligible contract, map each local update history to its current component tuple, active monitor memories and unexecuted commands. Histories with the same image admit the same next events and all the same local outcomes, agree on safety, and have the same handover targets. This map gives a label- and outcome-preserving correspondence between local update histories and paths of the generated game. A history-dependent policy satisfying the source obligations exists exactly when a winning policy exists in that game.
\end{proposition}
\begin{proof}
Let $F(h)$ be the tuple obtained by folding the stated source operations over history $h$, with pending commands the complement of executed commands. At entry, $F(h)=\Init(z)$, giving exactly $Q_0$. Suppose a history ends at $(\xi,\mu,P)$. For an ordinary event, both descriptions require the same participating component edges, retain their full Cartesian product and advance every active monitor. For a transfer, both require quiescence, pending status, completed predecessors and an outcome of the same $g_i$; they change only component $i$, advance active monitors and consume that command. A stop removes its own monitor without consuming the boundary and advances the others. A start advances existing monitors and installs the same initializer, without a second observation. Thus each permitted label and each outcome has exactly the same successor image in both descriptions; no choice of outcome is introduced or removed. Quiescence counts unsafe UC outcomes in both. Safety and compatible targets depend on identical coordinates.

Induction gives a game path for every source history. Conversely choose an old tuple witnessing the first game state and lift each edge by the one-step property, retaining its outcome. The same recurrence lifts infinite paths. Histories with the same image admit identical future extensions because old control is released at entry. A game policy therefore lifts through $F$. Conversely, fix one old witness for each entry and consistently lift each game path from it; applying a source policy to that lifted history defines a game history-dependent policy. Both translations preserve every enabled UC event, every selected-label outcome, unsafe prefixes, deadlocks and infinite avoiding runs. Winning existence is therefore equivalent for history-dependent policies. Theorem~\ref{thm:correctness}'s losing-closure argument rules out every history-dependent winner on failure and returns a state-based winner on success, establishing the stated equivalence.
\end{proof}
This proves correspondence for the supplied contract; requirement meaning and Java frontend correctness remain separate obligations. Technical Appendix~\ref{ta:source_synthesis} gives the full source recurrences and expanded proof.
